\documentclass{article}
\usepackage[utf8]{inputenc}
\usepackage[letterpaper,margin=0.9in]{geometry}
\usepackage{amsmath,amssymb,amsthm}
\usepackage{mathtools}
\usepackage{booktabs}
\usepackage{multirow}
\usepackage{float}
\usepackage{graphicx}
\usepackage{algorithm}
\usepackage{algpseudocode}
\usepackage{hyperref}
\usepackage{xcolor}
\usepackage{bm}

% ----------------------------------------------------------------------
\newtheorem{theorem}{Theorem}
\newtheorem{proposition}{Proposition}
\newtheorem{lemma}{Lemma}
\newtheorem{corollary}{Corollary}
\newtheorem{assumption}{Assumption}
\newtheorem{definition}{Definition}
\newtheorem{remark}{Remark}

\DeclareMathOperator{\KL}{KL}
\DeclareMathOperator{\Var}{Var}
\DeclareMathOperator*{\argmin}{arg\,min}
\DeclareMathOperator*{\argmax}{arg\,max}
\newcommand{\E}{\mathbb{E}}
\newcommand{\R}{\mathbb{R}}
\newcommand{\X}{\mathcal{X}}
\newcommand{\Q}{\mathcal{Q}}
\newcommand{\eps}{\varepsilon}
\newcommand{\lstar}{\lambda^{\star}}
\newcommand{\qlam}{p^{\lambda}(\cdot\mid x_{t_k})}
\newcommand{\todo}[1]{{\color{red}[#1]}}

\title{Soft-Constrained Discrete Diffusion via Lagrangian Duality:\\ Adaptive Dual Guidance as Per-Step Information Projection}
\author{Anonymous}
\date{}

\begin{document}
\maketitle

\begin{abstract}
We study \emph{constrained generation} from a pretrained discrete diffusion model: sampling sequences (molecules, text) that satisfy an attribute constraint while staying close to the model's data distribution. Existing constrained discrete-diffusion samplers take one of two extremes. \emph{Hard-projection} methods enforce the constraint exactly at every reverse step, achieving zero violation but, on our QM9 benchmarks, collapsing validity and diversity as the feasible set is squeezed. \emph{Classifier guidance} instead tilts each step softly by a constraint-satisfaction score raised to a fixed global weight $\gamma$; this weight is a hand-tuned hyperparameter that exposes a sharp over/under-guidance dilemma. We unify and interpolate between these extremes by formulating constrained generation as a per-step \emph{constrained information projection}: at each reverse step we seek the kernel closest in KL to the unconditional model dynamics subject to an explicit budget on the expected constraint cost, $\E_q[-\log r_{t_{k+1}}(x)]\le C$, where $r_{t_k}(x)=\mathbb P_\theta[X_{t_K}\in S\mid X_{t_k}=x]\in[0,1]$ is the constraint-satisfaction score---the discriminator's probability that the trajectory ultimately lands in the feasible set $S$. We show (i) the projection has a closed form, the exponential tilt $q_\lambda\propto p_\theta\cdot r_{t_{k+1}}(x)^\lambda$, so classifier guidance is recovered as the Lagrangian-optimal projection with $\gamma=\lambda$ equal to the constraint's multiplier; (ii) hard projection is the zero-budget limit $C\to0$ ($\lambda\to\infty$), so the cost budget $C$ is a single dial that interpolates from the unconstrained model to hard projection; and (iii) because the tilted family is an exponential family in $\lambda$, the dual is smooth and concave with gradient and Hessian given in closed form by the mean and variance of the constraint log-score $\log r_{t_{k+1}}(x)$. We exploit (iii) to \emph{solve} for the multiplier online: \emph{Adaptive Dual Guidance} (first-order, one classifier evaluation per step) and \emph{Adaptive Dual Newton Guidance} (second-order, with mode tracking) adapt $\lambda$ per step and per sample to meet the budget with least distortion. On constrained molecular generation with QM9 (a synthetic-accessibility constraint), the adaptive budget escapes both the fixed-$\gamma$ dilemma and the hard-projection validity collapse: it reaches a constraint-violation level that fixed guidance attains only by destroying validity and diversity, while keeping validity and uniqueness in the lightly-constrained regime.
\end{abstract}

\section{Introduction}
Discrete diffusion models have become a leading class of generative models for inherently discrete data---molecules written as SMILES strings, natural language, and biological sequences---where they match or surpass autoregressive models while generating by parallel, iterative refinement of a fully noised sequence \cite{austin2021d3pm,lou2024sedd,sahoo2024mdlm}. Many of their highest-value uses are not \emph{unconditional} generation but \emph{constrained} generation: we want samples that exhibit a target attribute---a molecule that is synthesizable or novel, a sentence that is non-toxic---while remaining close to the learned data distribution so that they stay valid and diverse. The central question is how to steer a \emph{pretrained} model toward the attribute without retraining and without sacrificing the sample quality that makes the base model useful.

Two families of constrained discrete-diffusion samplers address this, at opposite extremes. \emph{Hard-projection} methods enforce feasibility exactly: at every reverse step they project the model's distribution onto the constraint set, so the constraint is satisfied by construction \cite{cardei2025cdd}. This guarantees zero violation, but when the feasible region is small relative to the model's mass, repeatedly projecting drives samples off the data manifold; on our QM9 benchmarks validity and diversity collapse. \emph{Classifier guidance} instead steers softly: it tilts each reverse step by a constraint-satisfaction score raised to a guidance weight $\gamma$, preserving more of the base dynamics \cite{dhariwal2021guidance,nisonoff2024discreteguidance}. But $\gamma$ is a single global hyperparameter, hand-tuned and without a principled meaning, and it exposes a sharp \emph{over/under-guidance dilemma}: too small and the attribute constraint is left unmet, too large and the samples leave the manifold and validity collapses. As we show on QM9, no single fixed $\gamma$ achieves both constraint satisfaction and high quality.

We show that this guidance weight is not arbitrary. We formulate per-step guidance as a \emph{constrained information projection}: at each reverse step, find the step distribution closest in Kullback--Leibler divergence to the pretrained kernel subject to an explicit budget $C$ on the expected constraint cost, $\E_q[-\log r_{t_{k+1}}(x)]\le C$, with $r_{t_k}(x)=\mathbb P_\theta[X_{t_K}\in S\mid X_{t_k}=x]$ the constraint-satisfaction score---the discriminator's probability that generation ultimately lands in the feasible set $S$. This single, interpretable budget reorganizes the whole picture. First, the projection has a closed form---the exponential tilt $q_\lambda\propto p_\theta\cdot r_{t_{k+1}}(x)^\lambda$---so classifier guidance is recovered \emph{exactly} as the Lagrangian-optimal projection, with the guidance weight $\gamma=\lambda$ equal to the Lagrange multiplier of the constraint. The principled object to control is therefore not the weight $\gamma$ but the budget $C$, from which the correct $\lambda$ is \emph{solved}, not tuned. Second, hard projection is the zero-budget limit $C\!\to\!0$ ($\lambda\!\to\!\infty$): the two extreme families are two ends of one axis, continuously dialed by $C$. Third, because the tilted family is an exponential family in $\lambda$, the dual objective is smooth and concave with gradient and Hessian given in closed form by the \emph{mean} and \emph{variance} of the constraint log-score $\log r_{t_{k+1}}(x)$---quantities already available during guided decoding.

We exploit this structure to \emph{solve} for the multiplier online, adapting it per step and per sample so as to meet the budget with the least distortion to the base model. \emph{Adaptive Dual Guidance} takes one projected-gradient-ascent step on the dual per reverse step, at the same cost as standard guidance (one classifier evaluation), and is a drop-in replacement for fixed-$\gamma$ guidance. \emph{Adaptive Dual Newton Guidance} uses the closed-form curvature for second-order updates and tracks the running mode for a MAP-style output. On QM9 synthetic-accessibility-constrained generation, the adaptive budget escapes both failure modes: it reaches a constraint-violation level that fixed guidance attains only by collapsing validity and diversity, while holding validity and uniqueness in the lightly-constrained regime.

\paragraph{Contributions.}
\begin{itemize}
\item \textbf{Reframing.} We cast per-step classifier guidance for discrete diffusion as a constrained information projection, and identify the guidance weight with the Lagrange multiplier of an explicit expected-surprisal budget $C$ (Prop.~\ref{prop:tilt}).
\item \textbf{Unification.} We prove hard-projection constrained diffusion (CDD) is the zero-budget limit $C\!\to\!0$, $\lambda\!\to\!\infty$ of the same family (Prop.~\ref{prop:hard}), placing hard and soft constrained diffusion on a single axis controlled by $C$.
\item \textbf{Solver.} We expose the exponential-family structure of the dual---concave, with closed-form mean/variance gradient and Hessian (Prop.~\ref{prop:dual})---and derive two online algorithms (first- and second-order) with convergence and asymptotic-feasibility guarantees (Prop.~\ref{prop:converge}).
\item \textbf{Empirics.} On QM9 SA-constrained generation, adaptive $\lambda$ escapes the fixed-$\gamma$ over/under-guidance dilemma and the hard-projection validity collapse, dominating the fixed-$\gamma$ sweep on violation, validity, and diversity simultaneously.
\end{itemize}

\section{Background: Constrained Discrete Diffusion}
\label{sec:background}

\paragraph{Discrete diffusion.} A discrete diffusion model is defined over length-$L$ sequences $\X^L$ from a vocabulary $\X$ (with an absorbing \texttt{[MASK]} symbol adjoined for masked/absorbing-state models): a forward process noises a clean sequence and the model learns to reverse it. Sampling runs a learned \emph{unconditional} reverse process: along a time grid $0{=}t_0<t_1<\dots<t_K{=}T$, at each step $k$ it draws $x_{t_{k+1}}\sim p_\theta(\cdot\mid x_{t_k})$ from a kernel over $\X^L$, terminating in a clean sample $x_{t_K}$ distributed approximately as the data distribution $p_{\mathrm{data}}$. We write $X_{t_0},\dots,X_{t_K}$ for the resulting random trajectory (with realizations $x_{t_k}\in\X^L$), so $X_{t_K}$ is the clean terminal sample, and $\mathbb P_\theta,\E_\theta$ for probability and expectation under the law of this full reverse process---the composition of the one-step kernels $p_\theta(\cdot\mid x_{t_j})$ over $t_j\ge t_k$. We treat the pretrained one-step kernel $p_\theta(\cdot\mid x_{t_k})$ as a fixed primitive.

\paragraph{The constrained generation problem.} In constrained discrete diffusion the constraint is encoded by a \emph{feasible set} $S\subseteq\X^L$ of clean samples that satisfy it. The natural score of a partially-denoised state $x$ at step $t_k$ is the probability that the reverse process, continued from $x$, \emph{ultimately lands in} $S$,
\begin{equation}
\label{eq:rscore}
r_{t_k}(x)\;\coloneqq\;\E_\theta\!\big[\mathbf 1\{X_{t_K}\in S\}\,\big|\,X_{t_k}=x\big]
\;=\;\mathbb P_\theta\!\big[\,X_{t_K}\in S \,\big|\, X_{t_k}=x\,\big]\;\in\;[0,1].
\end{equation}
This is exactly what a time-conditioned \emph{classifier} (\emph{discriminator}) is trained to predict from noised states (D-CBG): how likely the constraint is to hold at the end of generation, sharpening from a soft score early in decoding to the indicator $r_{t_K}(x)=\mathbf 1[x\in S]$ at the terminal step. Its negative log $-\log r_{t_k}(x)\ge0$ is the \emph{constraint cost} (surprisal) of state $x$; for the soft method we take a strictly positive estimate $r_{t_k}\in(0,1]$, so this cost is finite. We seek samples that ultimately land in $S$ (low constraint cost) while staying close to $p_{\mathrm{data}}$. Two regimes have been studied:
\begin{itemize}
\item \emph{Hard / feasibility constraints}: require $X_{t_K}\in S$. Methods such as Constrained Discrete Diffusion (CDD) \cite{cardei2025cdd} enforce $S$ by \emph{projecting} the per-step distribution onto the feasible set, achieving $0\%$ violation by construction.
\item \emph{Soft / expectation constraints}: bound the \emph{expected} per-step constraint cost, $\E[-\log r_{t_{k+1}}(X_{t_{k+1}})]\le C$, trading a controlled amount of violation for higher quality.
\end{itemize}
Hard projection is uncompromising on violation but, when the feasible set is small relative to the model's mass, repeatedly projecting drives samples off the valid-data manifold---in our experiments validity and diversity collapse (Sec.~\ref{sec:experiments}). This motivates a principled \emph{soft} constraint with a tunable budget.

\paragraph{Estimating the constraint score.} The score \eqref{eq:rscore} is intractable in closed form, so we learn a time-conditioned discriminator $r_\phi(x,t_k)\in(0,1)$ to approximate it (D-CBG), fit by binary cross-entropy against the terminal feasibility label $y(x)=\mathbf 1[x\in S]$. The training pairs $(x_{t_k},t_k,y)$ can be drawn from either of two couplings between a clean molecule and its noised state.

\emph{(i) Data $+$ forward process.} Draw a clean datum from the data distribution, $x_{t_K}\sim p_{\mathrm{data}}$, label it by its own feasibility $y=\mathbf 1[x_{t_K}\in S]$, and noise it with the forward kernel $q(x_{t_k}\mid x_{t_K})$ to a randomly sampled level $t_k$:
\begin{equation}
\label{eq:cls-data}
\mathcal L_{\mathrm{data}}(\phi)=-\,\E_{x_{t_K}\sim p_{\mathrm{data}}}\;\E_{k}\;\E_{x_{t_k}\sim q(\cdot\mid x_{t_K})}\Big[\,y\log r_\phi(x_{t_k},t_k)+(1-y)\log\big(1-r_\phi(x_{t_k},t_k)\big)\Big].
\end{equation}

\emph{(ii) Pretrained-model path.} Draw a full unconditional trajectory from the pretrained reverse process, $(x_{t_0},\dots,x_{t_K})\sim p_\theta$, and label \emph{every} state along it by the trajectory's own terminal outcome $y=\mathbf 1[x_{t_K}\in S]$:
\begin{equation}
\label{eq:cls-model}
\mathcal L_{\mathrm{model}}(\phi)=-\,\E_{(x_{t_0},\dots,x_{t_K})\sim p_\theta}\;\E_{k}\Big[\,y\log r_\phi(x_{t_k},t_k)+(1-y)\log\big(1-r_\phi(x_{t_k},t_k)\big)\Big].
\end{equation}

The Bayes-optimal minimizer of \eqref{eq:cls-model} is \emph{exactly} the score \eqref{eq:rscore}, $r_\phi(x,t_k)\to\mathbb P_\theta[X_{t_K}\in S\mid X_{t_k}=x]$, because the labels are generated under the model's own reverse law---the route scores the very trajectories guidance will steer, removing any train/inference mismatch but requiring samples from the base model. The data route \eqref{eq:cls-data} instead targets the posterior $\mathbb P[x_{t_K}\in S\mid x_{t_k}=x]$ under the data-and-forward coupling; it needs only labelled clean molecules and the cheap forward kernel, and coincides with \eqref{eq:rscore} when the pretrained reverse process matches the forward posterior (a well-trained base model), differing otherwise.

\paragraph{Guidance as a soft heuristic.} The standard soft mechanism, classifier-based guidance (D-CBG; the constraint score $r_{t_k}$ is classically a classifier likelihood, whence the name), replaces the unconditional kernel by the \emph{tilted} step
\begin{equation}
\label{eq:cg-step}
p^{\gamma}(x\mid x_{t_k})\;\propto\;p_\theta(x\mid x_{t_k})\,r_{t_{k+1}}(x)^{\gamma},
\qquad \gamma\ge 0,
\end{equation}
sampling $x_{t_{k+1}}\sim p^\gamma(\cdot\mid x_{t_k})$. The guidance weight $\gamma$ is a single global hyperparameter shared across all steps and all samples; it is tuned by hand and, as we show, exposes a sharp over/under-guidance dilemma (small $\gamma$: constraint unmet; large $\gamma$: off-manifold collapse). Section~\ref{sec:method} shows that \eqref{eq:cg-step} is not a heuristic at all but the exact solution of a per-step soft-constrained projection, that hard projection is its $C\to0$ limit, and that the proper object to control is the budget $C$---from which the right $\gamma=\lambda$ is \emph{solved}, not tuned.

% ======================================================================
\section{Constrained Discrete Diffusion as Per-Step Information Projection}
\label{sec:method}

We develop the method around a single reverse step. The idea is to cast guided sampling as the information projection of the model's own step onto an explicit constraint-cost budget (Sec.~\ref{sec:perstep}), and to draw three consequences. The projection has a closed form with the functional form of classifier guidance, the guidance weight playing the role of a Lagrange multiplier (Prop.~\ref{prop:tilt}); the hard-projection regime is the zero-budget limit---exactly hard projection onto the feasible set at the terminal step---placing the two known regimes as endpoints of one continuum (Prop.~\ref{prop:hard}); and the multiplier maximizes a smooth concave dual whose gradient and Hessian are the closed-form mean and variance of the constraint score (Prop.~\ref{prop:dual}), which Sec.~\ref{sec:algorithms} turns into a cheap online solver. All statements are local to the step from $x_{t_k}$ and rest on Assumption~\ref{ass:standing} (finite support, strictly positive score, strictly feasible budget).

\subsection{The per-step constrained projection}
\label{sec:perstep}
Fix a reverse step from state $x_{t_k}$ and recall the pretrained base kernel $p_\theta(\cdot\mid x_{t_k})$ over $\X^L$ (Sec.~\ref{sec:background}). Constrained guidance replaces it by a \emph{guided} step distribution $q(\cdot\mid x_{t_k})$ that stays as close as possible to $p_\theta(\cdot\mid x_{t_k})$ while being constraint-aware to a controlled degree: ``closeness'' is measured by $\KL\big(q(\cdot\mid x_{t_k})\,\Vert\,p_\theta(\cdot\mid x_{t_k})\big)$ and ``constraint-awareness'' by an upper bound on the expected constraint cost $-\log r_{t_{k+1}}$ under the guided step. With $q(\cdot\mid x_{t_k})$ ranging over the probability simplex $\Q=\{q:q(x)\ge0,\ \sum_x q(x)=1\}$ on $\X^L$, this yields the per-step primal problem
\begin{equation}
\label{eq:primal}
\min_{q(\cdot\mid x_{t_k})\in\Q}\; \KL\big(q(\cdot\mid x_{t_k})\,\big\Vert\, p_\theta(\cdot\mid x_{t_k})\big)
\qquad\text{s.t.}\qquad
\E_{q(\cdot\mid x_{t_k})}\!\left[-\log r_{t_{k+1}}(x)\right]\;\le\;C,
\end{equation}
where $C\ge 0$ is a surprisal budget and the next state $x_{t_{k+1}}$ is drawn from the optimal $q(\cdot\mid x_{t_k})$. The constraint indicates: under the guided step, the constraint score $r_{t_{k+1}}(x)$ should be \emph{at least} $e^{-C}$ in expectation (equivalently, $\E_{q(\cdot\mid x_{t_k})}[\log r_{t_{k+1}}(x)]\ge -C$). Small $C$ demands an aggressively constraint-aware step; large $C$ recovers the base kernel $p_\theta(\cdot\mid x_{t_k})$.

\begin{assumption}[Standing assumptions for the per-step projection]
\label{ass:standing}
We analyze \eqref{eq:primal} on the support $\X_p\coloneqq\{x\in\X^L:p_\theta(x\mid x_{t_k})>0\}$ (finite, as $\X^L$ is) over an \emph{ideal joint} next-state distribution on $\X_p$; the practical sampler realizes this only approximately, within its (factorized) reverse-kernel family (Sec.~\ref{sec:algorithms}). The score is strictly positive on the support, $r_{t_{k+1}}\in(0,1]$ on $\X_p$, so $-\log r_{t_{k+1}}$ is finite there; write $m^\star\coloneqq\max_{x\in\X_p}r_{t_{k+1}}(x)$ and $c_{\min}\coloneqq-\log m^\star=\min_{x\in\X_p}\!\big(-\log r_{t_{k+1}}(x)\big)$, the smallest achievable cost. We restrict $\Q$ to distributions supported on $\X_p$, so $\KL(q\Vert p_\theta)<\infty$. A budget is \emph{strictly feasible} (Slater) iff $C>c_{\min}$; $C=c_{\min}$ is the boundary (Prop.~\ref{prop:hard}) and $C<c_{\min}$ is infeasible.
\end{assumption}

\begin{remark}[Exact conditioning at $\lambda{=}1$; the per-step budget is a local surrogate]
\label{rmk:myopic}
The score $r_{t_k}(x)=\mathbb P_\theta[X_{t_K}\in S\mid X_{t_k}=x]$ is the \emph{lookahead value} of conditioning on $\{X_{t_K}\in S\}$. By the tower property under the base process, $r_{t_k}(x_{t_k})=\E_{X_{t_{k+1}}\sim p_\theta(\cdot\mid x_{t_k})}[r_{t_{k+1}}(X_{t_{k+1}})]$, so $\{r_{t_k}(X_{t_k})\}_k$ is a $\mathbb P_\theta$-martingale and the exact Doob $h$-transform that samples the process conditioned on $\{X_{t_K}\in S\}$ is the \emph{first-power} tilt $p^{1}(x\mid x_{t_k})=p_\theta(x\mid x_{t_k})\,r_{t_{k+1}}(x)/r_{t_k}(x_{t_k})$ (its normalizer is exactly $r_{t_k}(x_{t_k})$). Our family $p^{\lambda}$ recovers this exact conditioning \emph{only} at $\lambda=1$; for $\lambda\neq1$ it is a power-tilt that trades exact conditioning for the soft budget $C$, with $\lambda$ adapting to $C$ per step. (The score $r$ is intractable in closed form, integrating over all future denoising \cite{wu2023practical,zhao2024probabilistic}, and is estimated by a time-conditioned discriminator, D-CBG; our analysis is orthogonal to that estimate.) Finally, the per-step constraint is a \emph{local surrogate} for the terminal constraint, not a proven relaxation of it: by Jensen it guarantees only the one-step bound $\E_{q}[r_{t_{k+1}}]\ge e^{-C}$, and controlling terminal feasibility under the \emph{guided} process would require a separate telescoping argument that we do not claim.
\end{remark}

The geometry of \eqref{eq:primal} is benign.

\begin{lemma}[Convexity]
\label{lem:convex}
The objective $q(\cdot\mid x_{t_k})\mapsto\KL(q(\cdot\mid x_{t_k})\Vert p_\theta(\cdot\mid x_{t_k}))$ is strictly convex on $\Q$, and the feasible set
$\mathcal{F}_C=\{q(\cdot\mid x_{t_k})\in\Q:\E_{q(\cdot\mid x_{t_k})}[-\log r_{t_{k+1}}(x)]\le C\}$ is convex and closed. Hence \eqref{eq:primal} is a convex program with a unique minimizer whenever $\mathcal{F}_C\neq\emptyset$.
\end{lemma}
\begin{proof}
\begin{align*}
\KL\big(q(\cdot\mid x_{t_k})\,\Vert\,p_\theta(\cdot\mid x_{t_k})\big)
  &=\sum_x q(x\mid x_{t_k})\log\frac{q(x\mid x_{t_k})}{p_\theta(x\mid x_{t_k})}
   \quad\text{strictly convex in }q\ (u\mapsto u\log u),\\
\E_{q(\cdot\mid x_{t_k})}[-\log r_{t_{k+1}}(x)]
  &=\sum_x q(x\mid x_{t_k})\,\big(-\log r_{t_{k+1}}(x)\big)
   \quad\text{linear in }q,\\
\mathcal{F}_C
  &=\Q\cap\big\{q(\cdot\mid x_{t_k}):\E_{q(\cdot\mid x_{t_k})}[-\log r_{t_{k+1}}]\le C\big\}
   \quad\text{(convex, closed)}.
\end{align*}
A strictly convex function attains a unique minimizer on the nonempty compact convex set $\mathcal{F}_C$.
\end{proof}

\subsection{The optimal projection is classifier guidance}
We now solve \eqref{eq:primal} in closed form via Lagrangian duality and show that the solution is exactly the tilted step \eqref{eq:cg-step}.

\begin{proposition}[Optimal tilted projection]
\label{prop:tilt}
Under Assumption~\ref{ass:standing} with a strictly feasible budget $C>c_{\min}$ (Slater), \eqref{eq:primal} admits a finite optimal multiplier $\lstar\in[0,\infty)$ and the unique solution is
\begin{equation}
\label{eq:qlam}
q^{*}(x\mid x_{t_k})\;=\;p^{\lstar}(x\mid x_{t_k})\;=\;\frac{p_\theta(x\mid x_{t_k})\,r_{t_{k+1}}(x)^{\lstar}}{Z(\lstar)},
\qquad
Z(\lambda)\coloneqq\sum_{x\in\X^L}p_\theta(x\mid x_{t_k})\,r_{t_{k+1}}(x)^{\lambda},
\end{equation}
for a multiplier $\lstar\ge 0$, where $\lstar=0$ iff the unconditional kernel already satisfies the constraint ($\E_{p_\theta(\cdot\mid x_{t_k})}[-\log r_{t_{k+1}}(x)]\le C$), and otherwise $\lstar>0$ is chosen so the constraint is active, $\E_{p^{\lstar}(\cdot\mid x_{t_k})}[-\log r_{t_{k+1}}(x)]=C$. In particular $p^{\lstar}(\cdot\mid x_{t_k})$ has the functional form of classifier guidance \eqref{eq:cg-step} with $\gamma=\lstar$.
\end{proposition}

\begin{proof}
Introduce a multiplier $\lambda\ge 0$ for the inequality constraint. The Lagrangian is
\begin{equation}
\label{eq:lagrangian}
\begin{aligned}
\mathcal{L}(q(\cdot\mid x_{t_k}),\lambda)
&=\KL(q(\cdot\mid x_{t_k})\Vert p_\theta(\cdot\mid x_{t_k}))+\lambda\!\left(\E_{q(\cdot\mid x_{t_k})}[-\log r_{t_{k+1}}(x)]-C\right)\\
&=\sum_x q(x\mid x_{t_k})\log\frac{q(x\mid x_{t_k})}{p_\theta(x\mid x_{t_k})}
-\lambda\sum_x q(x\mid x_{t_k})\log r_{t_{k+1}}(x)-\lambda C.
\end{aligned}
\end{equation}
Minimize over $q(\cdot\mid x_{t_k})\in\Q$ with a normalization multiplier $\nu$ for $\sum_x q(x\mid x_{t_k})=1$:
\begin{align*}
0&=\frac{\partial}{\partial q(x\mid x_{t_k})}\Big[\mathcal L+\nu\big(\textstyle\sum_x q(x\mid x_{t_k})-1\big)\Big]
  =\log\frac{q(x\mid x_{t_k})}{p_\theta(x\mid x_{t_k})}+1-\lambda\log r_{t_{k+1}}(x)+\nu,\\
\Longrightarrow\quad q(x\mid x_{t_k})&\propto p_\theta(x\mid x_{t_k})\,r_{t_{k+1}}(x)^{\lambda}=p^{\lambda}(x\mid x_{t_k})
  \quad\text{(unique inner minimizer, Lemma~\ref{lem:convex})}.
\end{align*}
Since $C>c_{\min}$, strong duality holds with zero gap and a finite maximizer $\lstar<\infty$ (Prop.~\ref{prop:strongduality}), so the primal optimum is $p^{\lstar}(\cdot\mid x_{t_k})$. Complementary slackness gives
\begin{equation*}
\lstar\big(\E_{p^{\lstar}(\cdot\mid x_{t_k})}[-\log r_{t_{k+1}}]-C\big)=0
\ \Longrightarrow\
\begin{cases}
\lstar=0: & p^{\lstar}(\cdot\mid x_{t_k})=p_\theta(\cdot\mid x_{t_k})\ \ (\text{if }p_\theta\text{ already feasible}),\\[2pt]
\lstar>0: & \E_{p^{\lstar}(\cdot\mid x_{t_k})}[-\log r_{t_{k+1}}]=C.
\end{cases}
\end{equation*}
\end{proof}

Proposition~\ref{prop:tilt} is the conceptual core of the paper. It says that the tilted step \eqref{eq:cg-step} \emph{is} the information projection of the model's own step onto the constraint set $\mathcal F_C$, with the guidance weight playing the role of the constraint's Lagrange multiplier. Classifier guidance is thus not an ad hoc reweighting but the functional form of this optimal projection; conversely, a fixed weight $\gamma$ is the optimal projection only for its own \emph{implicit} budget $C(\gamma)=\E_{p^{\gamma}(\cdot\mid x_{t_k})}[-\log r_{t_{k+1}}(x)]$, which varies from step to step and sample to sample. Holding $\gamma$ global therefore pins the implicit budget nowhere in particular---tight where feasibility is easy, slack where it is hard---which is exactly what creates the over/under-guidance dilemma. The principled object is not $\gamma$ but the budget $C$, from which the constraint-satisfying $\lambda$ is solved per step.

\subsection{Hard projection is the zero-budget limit}
The same projection connects the soft and hard regimes of Sec.~\ref{sec:background}: shrinking the budget $C$ to its minimum $c_{\min}$ drives the multiplier to infinity and collapses each step onto the score-maximizing states---which, at the terminal transition, is exactly hard projection onto the feasible set $S$.

\begin{proposition}[Hard projection as $C\to0$ / $\lambda\to\infty$]
\label{prop:hard}
Let $\X_p=\{x\in\X^L:p_\theta(x\mid x_{t_k})>0\}$ and $m^\star=\max_{x\in\X_p}r_{t_{k+1}}(x)$ with maximizer set $\X^\star=\{x\in\X_p:r_{t_{k+1}}(x)=m^\star\}$. As $\lambda\to\infty$, the tilt $p^{\lambda}(\cdot\mid x_{t_k})$ of \eqref{eq:qlam} converges to the base kernel restricted to the score-maximizing states,
\begin{equation}
p^{\lambda}(x\mid x_{t_k})\ \xrightarrow{\lambda\to\infty}\
p^{\infty}(x\mid x_{t_k})=\frac{p_\theta(x\mid x_{t_k})\,\mathbf 1[x\in\X^\star]}{\sum_{x'\in\X^\star}p_\theta(x'\mid x_{t_k})},
\end{equation}
and $\E_{p^{\lambda}(\cdot\mid x_{t_k})}[-\log r_{t_{k+1}}(x)]\downarrow c_{\min}=-\log m^\star$, the smallest achievable cost. Hence \eqref{eq:primal} is feasible iff $C\ge c_{\min}$; at the boundary $C=c_{\min}$ the optimum is the limit $p^{\infty}(\cdot\mid x_{t_k})$ (the base kernel restricted to $\X^\star$ and renormalized), attained only as $\lstar\to\infty$. At the terminal transition $r_{t_K}=\mathbf 1[\cdot\in S]$, so $\X^\star=S\cap\X_p$ and $p^{\infty}$ is literally the hard projection onto the feasible set $S$; at earlier steps $\X^\star=\argmax_x r_{t_{k+1}}(x)$ is the set of most-feasible-leading next states, which need not equal $S$.
\end{proposition}
\begin{proof}
Set $\tilde r_{t_{k+1}}(x)\coloneqq r_{t_{k+1}}(x)/m^\star\le1$ and let $\lambda\to\infty$:
\begin{align*}
p^{\lambda}(x\mid x_{t_k})&\propto p_\theta(x\mid x_{t_k})\,\tilde r_{t_{k+1}}(x)^{\lambda},
  \qquad \tilde r_{t_{k+1}}(x)^{\lambda}\to\mathbf 1[x\in\X^\star]\ \ (\tilde r_{t_{k+1}}<1\ \text{off}\ \X^\star,\ =1\ \text{on}\ \X^\star),\\
p^{\lambda}(\cdot\mid x_{t_k})&\longrightarrow p^{\infty}(\cdot\mid x_{t_k})=\frac{p_\theta(x\mid x_{t_k})\,\mathbf 1[x\in\X^\star]}{\sum_{x'\in\X^\star}p_\theta(x'\mid x_{t_k})},\\
\frac{d}{d\lambda}\,\E_{p^{\lambda}(\cdot\mid x_{t_k})}[-\log r_{t_{k+1}}]
  &=-\Var_{p^{\lambda}(\cdot\mid x_{t_k})}(\log r_{t_{k+1}})\le0
  \qquad(\text{by }\eqref{eq:expfam-deriv}),\\
\E_{p^{\lambda}(\cdot\mid x_{t_k})}[-\log r_{t_{k+1}}]&\;\downarrow\;c_{\min}=\min_{x\in\X_p}\big(-\log r_{t_{k+1}}(x)\big).
\end{align*}
Feasibility of \eqref{eq:primal} therefore requires $C\ge c_{\min}$.
\end{proof}

Proposition~\ref{prop:hard} places hard-projection sampling and soft guidance on a single axis dialed by $C$. As $C\downarrow c_{\min}$ the multiplier diverges and each step collapses onto $\X^\star=\argmax_x r_{t_{k+1}}(x)$, the next states judged most likely to reach $S$. We emphasize the scope of the match with hard-projection methods such as CDD \cite{cardei2025cdd}: at the \emph{terminal} transition $\X^\star=S\cap\X_p$, so the zero-budget limit is exactly hard projection onto the feasible set $S$---CDD's set-projection primitive; at earlier steps it is the \emph{analogous} greedy restriction to the most-feasible-leading states rather than to $S$ itself, and a different hard sampler that projects onto a per-step relaxation of $S$ need not coincide with our limit there. Either way, discarding all off-$\X^\star$ mass at every step compounds over $K$ steps and sheds validity and diversity, which is why hard projection buys $0\%$ violation at the cost of collapse. A finite budget $C>c_{\min}$ instead keeps mass on near-optimal, on-manifold states, with $\lstar(C)<\infty$ measuring \emph{how hard} to project; our method sits strictly between the unconstrained model ($\lstar=0$) and this hard-projection endpoint ($\lstar\to\infty$), with $C$ the single dial.

\subsection{The dual problem and its exponential-family structure}
To turn this insight into an algorithm we study the dual function
\begin{equation}
\label{eq:dual}
g(\lambda)\;\coloneqq\;\min_{q(\cdot\mid x_{t_k})\in\Q}\mathcal L(q(\cdot\mid x_{t_k}),\lambda)
\;=\;-\log Z(\lambda)-\lambda C,
\end{equation}
where the second equality follows by substituting the inner minimizer $\qlam$ into \eqref{eq:lagrangian} (the KL plus cross-term telescopes to $-\log Z(\lambda)$). The \emph{dual problem} is
\begin{equation}
\label{eq:dual-problem}
\max_{\lambda\ge 0}\; g(\lambda).
\end{equation}

The key observation is that $\{\qlam\}_{\lambda\ge0}$ is a one-parameter \emph{exponential family} in the natural parameter $\lambda$, with base measure $p_\theta(\cdot\mid x_{t_k})$, sufficient statistic $T(x)=\log r_{t_{k+1}}(x)$, and log-partition $A(\lambda)=\log Z(\lambda)$:
\begin{equation}
p^{\lambda}(x\mid x_{t_k})=p_\theta(x\mid x_{t_k})\exp\!\big(\lambda\,T(x)-A(\lambda)\big).
\end{equation}
The cumulants of $T$ under $\qlam$ are therefore derivatives of the log-partition $A$, a standard fact we record as the workhorse of our solver.

\begin{lemma}[Exponential-family cumulant identities]
\label{lem:expfam}
Let $\{q_\theta\}_{\theta\in\R}$ be a one-parameter exponential family $q_\theta(x)=h(x)\exp\!\big(\theta\,T(x)-A(\theta)\big)$ on a finite set, with base measure $h\ge0$, sufficient statistic $T$, and log-partition $A(\theta)=\log\sum_x h(x)e^{\theta T(x)}$. Then $A$ is real-analytic and convex, and its first two derivatives are the mean and variance of $T$:
\begin{equation}
\label{eq:expfam-deriv}
A'(\theta)=\E_{q_\theta}\!\big[T(X)\big],
\qquad
A''(\theta)=\Var_{q_\theta}\!\big(T(X)\big)\;\ge\;0.
\end{equation}
\end{lemma}
\begin{proof}
On the finite set $\sum_x h(x)e^{\theta T(x)}\in(0,\infty)$ for every real $\theta$, so $A$ is real-analytic; differentiating,
\begin{align*}
A'(\theta)&=\frac{\sum_x h(x)\,T(x)\,e^{\theta T(x)}}{\sum_x h(x)\,e^{\theta T(x)}}=\E_{q_\theta}\big[T(X)\big],\\
A''(\theta)&=\E_{q_\theta}\big[T(X)^2\big]-\E_{q_\theta}\big[T(X)\big]^2=\Var_{q_\theta}\big(T(X)\big)\ge0.
\end{align*}
Nonnegativity of $A''$ gives convexity of $A$.
\end{proof}

\begin{proposition}[Smooth concave dual with mean/variance derivatives]
\label{prop:dual}
$g$ is real-analytic on $\R$ and concave on $\R_{\ge0}$, with
\begin{align}
g'(\lambda) &= -\,\E_{\qlam}\!\left[\log r_{t_{k+1}}(x)\right]-C
           \;=\;\E_{\qlam}\!\left[-\log r_{t_{k+1}}(x)\right]-C,
\label{eq:gprime}\\
g''(\lambda) &= -\,\Var_{\qlam}\!\left(\log r_{t_{k+1}}(x)\right)\;\le\;0.
\label{eq:gsecond}
\end{align}
Consequently $g'$ is nonincreasing. Under Assumption~\ref{ass:standing} with $C>c_{\min}$, a \emph{finite} maximizer $\lstar=\argmax_{\lambda\ge0}g(\lambda)<\infty$ is attained (Prop.~\ref{prop:strongduality}), unique whenever $\log r_{t_{k+1}}$ is non-degenerate under $p^{\lstar}(\cdot\mid x_{t_k})$ (so $g''(\lstar)<0$), and characterized by the complementarity
\begin{equation}
\lstar=0 \ \text{ and }\ g'(0)\le 0,
\qquad\text{or}\qquad
\lstar>0 \ \text{ and }\ g'(\lstar)=0\ \Leftrightarrow\ \E_{p^{\lstar}(\cdot\mid x_{t_k})}[-\log r_{t_{k+1}}(x)]=C.
\end{equation}
At the boundary $C=c_{\min}$ the supremum is approached only as $\lambda\to\infty$ (Prop.~\ref{prop:hard}); for $C<c_{\min}$, $g$ is unbounded above.
\end{proposition}
\begin{proof}
Apply Lemma~\ref{lem:expfam} with base measure $p_\theta(\cdot\mid x_{t_k})$, $T(x)=\log r_{t_{k+1}}(x)$, $A(\lambda)=\log Z(\lambda)$, together with $g(\lambda)=-A(\lambda)-\lambda C$ from \eqref{eq:dual}:
\begin{align*}
A'(\lambda)&=\E_{\qlam}[\log r_{t_{k+1}}],\qquad A''(\lambda)=\Var_{\qlam}(\log r_{t_{k+1}}),\\
g'(\lambda)&=-A'(\lambda)-C=\E_{\qlam}[-\log r_{t_{k+1}}]-C,\\
g''(\lambda)&=-A''(\lambda)=-\Var_{\qlam}(\log r_{t_{k+1}})\le0.
\end{align*}
Thus $g$ is concave; maximizing over $\R_{\ge0}$ yields the stated complementarity, and $g''<0$ gives uniqueness.
\end{proof}

\begin{proposition}[Strong duality and finite multiplier]
\label{prop:strongduality}
Under Assumption~\ref{ass:standing} with $C>c_{\min}$ (Slater), strong duality holds with zero gap, $\max_{\lambda\ge0}g(\lambda)=\min_{q(\cdot\mid x_{t_k})\in\mathcal F_C}\KL(q(\cdot\mid x_{t_k})\Vert p_\theta(\cdot\mid x_{t_k}))$; the dual maximizer $\lstar<\infty$ is attained, and the primal optimum is attained and unique (Lemma~\ref{lem:convex}).
\end{proposition}
\begin{proof}
The program \eqref{eq:primal} is convex (Lemma~\ref{lem:convex}) on the compact simplex over the finite $\X_p$, with $-\log r_{t_{k+1}}$ finite. \emph{Slater.} For any $x^\star\in\X^\star$,
\begin{align*}
\bar q&=(1-\delta)\,\delta_{x^\star}+\delta\,p_\theta(\cdot\mid x_{t_k})
  \quad\big(\text{supported on }\X_p,\ \KL(\bar q\Vert p_\theta(\cdot\mid x_{t_k}))<\infty\big),\\
\E_{\bar q}[-\log r_{t_{k+1}}]&\xrightarrow{\ \delta\to0\ }c_{\min}<C
  \quad\Longrightarrow\quad \bar q\ \text{is strictly feasible},
\end{align*}
so strong duality holds with zero gap and an attained dual maximizer (Boyd \& Vandenberghe, \S5.2.3). \emph{Finiteness of $\lstar$.} As $\lambda\to\infty$,
\begin{align*}
-\log Z(\lambda)&=\lambda\,c_{\min}-\log\!\Big(\sum_{x\in\X^\star}p_\theta(x\mid x_{t_k})\Big)+o(1),\\
g(\lambda)&=-\log Z(\lambda)-\lambda C=-\lambda\,(C-c_{\min})+O(1)\;\longrightarrow\;-\infty\quad(C>c_{\min}),
\end{align*}
so the maximizer $\lstar<\infty$. The primal optimum is unique by strict convexity on the compact feasible set (Lemma~\ref{lem:convex}).
\end{proof}

\paragraph{Why solve the dual.} The primal optimizes a distribution over the entire (exponentially large) state space $\X^L$, yet Prop.~\ref{prop:dual} shows the only unknown is the \emph{scalar} $\lstar$, and both $g'$ and $g''$ are expectations under $\qlam$---the very distribution we already sample from while decoding. A dual gradient therefore needs \emph{no model or classifier machinery beyond what guidance already evaluates}. Solving $g'(\lambda)=0$ exactly---enforcing $\E_{p^{\lambda}(\cdot\mid x_{t_k})}[\log r_{t_{k+1}}(x)]=-C$ by enumerating $\X^L$---is intractable at scale, so we instead ascend the dual with the cheap stochastic estimates of \eqref{eq:gprime}--\eqref{eq:gsecond} developed next.

% ======================================================================
\section{Algorithms}
\label{sec:algorithms}

\subsection{Adaptive Dual Guidance (first order)}
We solve the dual \eqref{eq:dual-problem} online, interleaving one projected-gradient-ascent step with each reverse diffusion step. Using the exact gradient $g'(\lambda_k)$ from \eqref{eq:gprime}, the \emph{theoretical} projected ascent step with size $\rho_k$ is
\begin{equation}
\label{eq:lambda-update-exact}
\lambda_{k+1}=\big(\lambda_k+\rho_k\,g'(\lambda_k)\big)_{+}
=\Big(\lambda_k+\rho_k\big(\E_{x\sim p^{\lambda_k}(\cdot\mid x_{t_k})}[-\log r_{t_{k+1}}(x)]-C\big)\Big)_{+}.
\end{equation}
The expectation in \eqref{eq:lambda-update-exact} is taken over the tilted next-step kernel $p^{\lambda_k}(\cdot\mid x_{t_k})$ at the current state $x_{t_k}$---\emph{exactly} the distribution from which we already draw $x_{t_{k+1}}$ to advance the chain. Rather than enumerate it, we reuse that single draw $x_{t_{k+1}}\sim p^{\lambda_k}(\cdot\mid x_{t_k})$ as an \emph{unbiased} one-sample estimate of the expectation, $\widehat{g'(\lambda_k)}=-\log r_{t_{k+1}}(x_{t_{k+1}})-C$, at no extra evaluation. Substituting into \eqref{eq:lambda-update-exact} gives the practical update
\begin{equation}
\label{eq:lambda-update}
\lambda_{k+1}=\Big(\lambda_k+\rho_k\,\widehat{g'(\lambda_k)}\Big)_{+}
=\Big(\lambda_k-\rho_k\big(\log r_{t_{k+1}}(x_{t_{k+1}})+C\big)\Big)_{+}.
\end{equation}
Here $(\cdot)_+=\max(\cdot,0)$ enforces $\lambda\ge0$. Intuitively: if the new state's constraint log-score $\log r_{t_{k+1}}(x_{t_{k+1}})$ is below the target $-C$ (constraint violated), the bracket is positive and $\lambda$ \emph{increases}, tightening guidance on the next step; once the constraint is comfortably satisfied, $\lambda$ relaxes back toward $0$ and distortion is released. The multiplier is maintained \emph{per sample}, so different trajectories in a batch receive different guidance strengths. In our implementation we additionally clip $\lambda$ to $[0,\lambda_{\max}]$ for stability. Each step costs exactly one classifier evaluation---the classifier is already evaluated to form the tilt, the same as standard guidance---so Adaptive Dual Guidance is a drop-in replacement for fixed-$\gamma$ D-CBG.

\begin{algorithm}[t]
\caption{Adaptive Dual Guidance}
\label{alg:adg}
\begin{algorithmic}[1]
\Require unconditional kernel $p_\theta(\cdot\mid x_{t_k})$; constraint scorer $r_{t_k}(\cdot)=\mathbb P_\theta[X_{t_K}\in S\mid X_{t_k}{=}\cdot]$; step sizes $\{\rho_k\}$; budget $C$; cap $\lambda_{\max}$; time grid $\{t_0,t_1,\dots,t_K\}$
\State initialize $\lambda_0$ (e.g.\ $0$), $x_{t_0}$
\For{$k=0,1,\dots,K-1$}
  \State $p^{\lambda_k}(\cdot\mid x_{t_k})\propto p_\theta(\cdot\mid x_{t_k})\,r_{t_{k+1}}(\cdot)^{\lambda_k}$ \Comment{tilted step, \eqref{eq:cg-step}}
  \State sample $x_{t_{k+1}}\sim p^{\lambda_k}(\cdot\mid x_{t_k})$
  \State $\lambda_{k+1}\gets \mathrm{clip}\big(\lambda_k-\rho_k(\log r_{t_{k+1}}(x_{t_{k+1}})+C),\,0,\,\lambda_{\max}\big)$ \Comment{dual ascent, \eqref{eq:lambda-update}}
\EndFor
\State \Return $x_{t_K}$
\end{algorithmic}
\end{algorithm}

\paragraph{Exact per-step solve (first order).}
Algorithm~\ref{alg:adg} takes a \emph{single} ascent step per reverse step: the updated multiplier is formed from the transition out of $x_{t_k}$ but applied only to the \emph{next} step, so $\lambda_k$ is a one-step-lagged tracker of the per-step optimum (Remark~\ref{rmk:onestep}). When this lag is undesirable---e.g.\ on coarse time grids, or when one wants the drawn transition to use exactly the projection-optimal multiplier of its \emph{own} step---we instead \emph{solve} the frozen-step dual $g_k$ before committing the draw. Because $x_{t_k}$, the base kernel $p_\theta(\cdot\mid x_{t_k})$, and the score $r_{t_{k+1}}$ are all fixed within a step, $g_k$ is a fixed scalar concave function of $\lambda$ (Prop.~\ref{prop:dual}); we ascend it by projected gradient ascent \eqref{eq:lambda-update}, repeatedly drawing from the tilted kernel $\qlam$ at the frozen state to estimate $g_k'(\lambda)=\E_{\qlam}[-\log r_{t_{k+1}}(x)]-C$ until $|\widehat{g_k'}|\le\eps_{\mathrm{tol}}$, then drawing the actual next state from the solved kernel $p^{\lstar_k}(\cdot\mid x_{t_k})$ (Algorithm~\ref{alg:inner-fo}). Dropping the variance estimate lets a single draw $n{=}1$ suffice per inner iteration, at the price of step-size tuning and, typically, more inner iterations $J$. The limits are transparent: with $J{=}1$ it reduces exactly to Algorithm~\ref{alg:adg}, and as $J\!\to\!\infty$ with a diminishing schedule $\{\rho_j\}$ it converges to the per-step optimum $\lstar_k$ by Proposition~\ref{prop:converge}(i),(iii).

\begin{algorithm}[t]
\caption{Adaptive Dual Guidance with exact per-step solve (first-order inner loop)}
\label{alg:inner-fo}
\begin{algorithmic}[1]
\Require unconditional kernel $p_\theta(\cdot\mid x_{t_k})$; constraint scorer $r_{t_k}(\cdot)$; budget $C$; cap $\lambda_{\max}$; inner iterations $J$, minibatch size $n\ge1$, inner step sizes $\{\rho_j\}$, tolerance $\eps_{\mathrm{tol}}$; time grid $\{t_0,\dots,t_K\}$
\State initialize $\lambda\gets\lambda_0$, $x_{t_0}$
\For{$k=0,1,\dots,K-1$}
  \State $\lambda'\gets\lambda$ \Comment{warm-start the inner solve from the previous step}
  \For{$j=1,\dots,J$} \Comment{solve the frozen-step dual $g_k$ at fixed $x_{t_k}$}
    \State sample $\{x^i\}_{i=1}^n\stackrel{\text{i.i.d.}}{\sim}p^{\lambda'}(\cdot\mid x_{t_k})\propto p_\theta(\cdot\mid x_{t_k})\,r_{t_{k+1}}(\cdot)^{\lambda'}$
    \State $\widehat{g_k'}\gets-\mathrm{Avg}_i\log r_{t_{k+1}}(x^i)-C$ \Comment{$\widehat{g_k'}$ unbiased for $g_k'(\lambda')$, \eqref{eq:gprime}}
    \If{$|\widehat{g_k'}|\le\eps_{\mathrm{tol}}$} \textbf{break} \EndIf
    \State $\lambda'\gets\mathrm{clip}\big(\lambda'+\rho_j\,\widehat{g_k'},\,0,\,\lambda_{\max}\big)$ \Comment{projected gradient ascent on $g_k$, \eqref{eq:lambda-update}}
  \EndFor
  \State $\lstar_k\gets\lambda'$ \Comment{(approx.) per-step optimal multiplier}
  \State sample $x_{t_{k+1}}\sim p^{\lstar_k}(\cdot\mid x_{t_k})$ \Comment{advance using \emph{this} step's solved multiplier}
  \State $\lambda\gets\lstar_k$
\EndFor
\State \Return $x_{t_K}$
\end{algorithmic}
\end{algorithm}

\subsection{Adaptive Dual Newton Guidance (second order)}
Because $g''(\lambda)=-\Var_{\qlam}(\log r_{t_{k+1}}(x))$ is available at the same cost as the gradient, we can take damped Newton steps on the concave dual, which removes the step-size sensitivity of \eqref{eq:lambda-update}. Given a minibatch $\{x^i\}_{i=1}^n\sim p^{\lambda_k}(\cdot\mid x_{t_k})$ we estimate
\begin{equation}
\widehat\mu=\frac1n\sum_{i=1}^n\log r_{t_{k+1}}(x^i),
\qquad
\widehat\sigma^2=\widehat{\Var}\big(\log r_{t_{k+1}}(x^i)\big),
\end{equation}
so that $\widehat{g'(\lambda_k)}=-\widehat\mu-C$ and $\widehat{g''(\lambda_k)}=-\widehat\sigma^2$. The projected Newton-ascent update $\lambda\!\leftarrow\!(\lambda-g'/g'')_+$ becomes
\begin{equation}
\label{eq:newton-update}
\lambda_{k+1}=\Big(\lambda_k-\frac{\widehat\mu+C}{\widehat\sigma^2}\Big)_{+}.
\end{equation}
Equation~\eqref{eq:newton-update} rescales the gradient by the local curvature: in low-variance (confident) regions it takes large steps, in high-variance regions it is cautious. We pair the second-order solver with \emph{mode tracking}: at each step we evaluate the tilted distribution at the running mode $x_{\mathrm{mode}}$ and return the mode rather than a random draw, yielding a MAP-style estimate of the constrained optimum. Algorithm~\ref{alg:newton} summarizes the procedure.

\begin{algorithm}[t]
\caption{Adaptive Dual Newton Guidance}
\label{alg:newton}
\begin{algorithmic}[1]
\Require as in Alg.~\ref{alg:adg}; sample size $n$
\State initialize $x_{t_0}$, $\lambda_0$, $\widehat\mu=\log r_{t_0}(x_{t_0})$, $\widehat\sigma^2=\sigma_0^2$, $x_{\mathrm{mode}}=x_{t_0}$
\For{$k=0,1,\dots,K-1$}
  \State $\lambda_{k+1}\gets\big(\lambda_k-(\widehat\mu+C)/\widehat\sigma^2\big)_+$ \Comment{Newton ascent, \eqref{eq:newton-update}}
  \State $p^{\lambda_{k+1}}(\cdot\mid x_{\mathrm{mode}})\propto p_\theta(\cdot\mid x_{\mathrm{mode}})\,r_{t_{k+1}}(\cdot)^{\lambda_{k+1}}$
  \State sample $\{x^i\}_{i=1}^n\stackrel{\text{i.i.d.}}{\sim}p^{\lambda_{k+1}}$
  \State $\widehat\mu\gets\mathrm{Avg}_i\log r_{t_{k+1}}(x^i)$;\quad $\widehat\sigma^2\gets\widehat{\Var}_i\log r_{t_{k+1}}(x^i)$;\quad $x_{\mathrm{mode}}\gets\mathrm{mode}(\{x^i\}_{i=1}^n)$
\EndFor
\State \Return $x_{t_K}=x_{\mathrm{mode}}$
\end{algorithmic}
\end{algorithm}

\paragraph{Exact per-step solve (second order).}
The same frozen-step solve admits a second-order inner loop, replacing the projected gradient step of Algorithm~\ref{alg:inner-fo} with the damped Newton step \eqref{eq:newton-update} that uses the closed-form curvature $\widehat{g_k''}=-\widehat\sigma^2$ alongside the gradient (Algorithm~\ref{alg:inner}). Because $g_k$ is a fixed scalar, strongly concave function of $\lambda$ with closed-form curvature (Lemma~\ref{lem:expfam}), Newton converges in a handful of inner steps; and because the per-step optima drift slowly along the trajectory (the duals $g_{k-1},g_k$ are close; Remark~\ref{rmk:onestep}), warm-starting $\lambda'$ from $\lstar_{k-1}$ keeps $J$ small---often $J{=}1$, recovering Algorithm~\ref{alg:newton} as the single-inner-step limit. Both inner-loop solvers (Algorithms~\ref{alg:inner-fo} and~\ref{alg:inner}) cost up to $J\cdot n$ classifier evaluations per reverse step versus one for the online schemes (Algorithms~\ref{alg:adg} and~\ref{alg:newton})---the price of removing the lag---and both realize the frozen-$k$ idealization of Proposition~\ref{prop:converge}(i)--(ii): driving $\widehat{g_k'}\to0$ enforces $\E_{p^{\lstar_k}(\cdot\mid x_{t_k})}[-\log r_{t_{k+1}}]\to C$ on each step's own transition $x_{t_k}\to x_{t_{k+1}}$.

\begin{algorithm}[t]
\caption{Adaptive Dual Guidance with exact per-step solve (inner loop)}
\label{alg:inner}
\begin{algorithmic}[1]
\Require unconditional kernel $p_\theta(\cdot\mid x_{t_k})$; constraint scorer $r_{t_k}(\cdot)$; budget $C$; cap $\lambda_{\max}$; inner iterations $J$, minibatch size $n$, tolerance $\eps_{\mathrm{tol}}$; time grid $\{t_0,\dots,t_K\}$
\State initialize $\lambda\gets\lambda_0$, $x_{t_0}$
\For{$k=0,1,\dots,K-1$}
  \State $\lambda'\gets\lambda$ \Comment{warm-start the inner solve from the previous step}
  \For{$j=1,\dots,J$} \Comment{solve the frozen-step dual $g_k$ at fixed $x_{t_k}$}
    \State sample $\{x^i\}_{i=1}^n\stackrel{\text{i.i.d.}}{\sim}p^{\lambda'}(\cdot\mid x_{t_k})\propto p_\theta(\cdot\mid x_{t_k})\,r_{t_{k+1}}(\cdot)^{\lambda'}$
    \State $\widehat\mu\gets\mathrm{Avg}_i\log r_{t_{k+1}}(x^i)$;\quad $\widehat\sigma^2\gets\widehat{\Var}_i\log r_{t_{k+1}}(x^i)$
    \State $\widehat{g_k'}\gets-\widehat\mu-C$ \Comment{$\widehat{g_k'}$ unbiased for $g_k'(\lambda')$, \eqref{eq:gprime}}
    \If{$|\widehat{g_k'}|\le\eps_{\mathrm{tol}}$} \textbf{break} \EndIf
    \State $\lambda'\gets\mathrm{clip}\big(\lambda'-\widehat{g_k'}/(-\widehat\sigma^2),\,0,\,\lambda_{\max}\big)$ \Comment{Newton step on $g_k$; or $(\lambda'+\rho\,\widehat{g_k'})_+$}
  \EndFor
  \State $\lstar_k\gets\lambda'$ \Comment{(approx.) per-step optimal multiplier}
  \State sample $x_{t_{k+1}}\sim p^{\lstar_k}(\cdot\mid x_{t_k})$ \Comment{advance using \emph{this} step's solved multiplier}
  \State $\lambda\gets\lstar_k$
\EndFor
\State \Return $x_{t_K}$
\end{algorithmic}
\end{algorithm}

\subsection{Convergence of the dual solver}
We record the standard guarantee that the inner dual ascent reaches the constraint-satisfying multiplier; this is what justifies calling $\lambda_k$ ``adaptive'' rather than merely ``time-varying''.

\begin{proposition}[Dual convergence and asymptotic feasibility]
\label{prop:converge}
Fix the kernel $p_\theta$ and budget $C$, and suppose Slater's condition (Prop.~\ref{prop:strongduality}) holds so that $\lstar<\infty$. Let $L=\sup_x|{-}\log r_{t_{k+1}}(x)-C|<\infty$ (finite on $\X^L$).
\begin{enumerate}
\item[(i)] \textbf{Exact gradients.} If $\lambda_{k+1}=(\lambda_k+\rho_k g'(\lambda_k))_+$ with $g''\le-m<0$ on the relevant interval and constant step $\rho_k=\rho\in(0,2/M)$ where $M=\sup|g''|$, then $\lambda_k\to\lstar$ linearly. With diminishing steps $\rho_k\to0$, $\sum_k\rho_k=\infty$, projected subgradient ascent satisfies $g(\lambda_k)\to g(\lstar)$, hence (strong duality) $\KL(p^{\lambda_k}(\cdot\mid x_{t_k})\Vert p_\theta(\cdot\mid x_{t_k}))\to\KL(p^{\lstar}(\cdot\mid x_{t_k})\Vert p_\theta(\cdot\mid x_{t_k}))$ and the constraint is satisfied in the limit.
\item[(ii)] \textbf{Newton.} On the interval where $g''\le -m<0$, the damped projected Newton iteration \eqref{eq:newton-update} with exact moments converges locally quadratically to $\lstar$.
\item[(iii)] \textbf{Stochastic gradients.} With the one-sample estimator of \eqref{eq:lambda-update}, $\E[\widehat{g'(\lambda_k)}\mid\lambda_k]=g'(\lambda_k)$ is unbiased and the estimator is bounded by $L$; projected stochastic gradient ascent with $\rho_k=\rho_0/\sqrt{k}$ gives $\E[g(\lstar)-g(\bar\lambda_K)]=O(L/\sqrt{K})$ for the running average $\bar\lambda_K$.
\end{enumerate}
\end{proposition}
\begin{proof}[Proof sketch]
$g$ is concave (Prop.~\ref{prop:dual}) and, on any bounded $\lambda$-interval, $L$-Lipschitz with $M$-Lipschitz gradient. Parts~(i),(iii) are the textbook projected (stochastic) gradient-ascent rates over $\R_{\ge0}$; the one-sample gradient is unbiased,
\begin{equation*}
\E\big[\widehat{g'(\lambda_k)}\,\big|\,\lambda_k\big]
=\E_{p^{\lambda_k}(\cdot\mid x_{t_k})}[-\log r_{t_{k+1}}]-C
=g'(\lambda_k)\quad(\text{by }\eqref{eq:gprime}),\qquad x_{t_{k+1}}\sim p^{\lambda_k}(\cdot\mid x_{t_k}).
\end{equation*}
Part~(ii) is Newton's local quadratic rate for a $C^2$, strongly concave objective with locally Lipschitz Hessian, preserved under projection. Asymptotic feasibility follows from strong duality and complementary slackness (Prop.~\ref{prop:strongduality}).
\end{proof}

\begin{remark}[One ascent step per reverse step]
\label{rmk:onestep}
In practice we do not run the dual solver to convergence at a frozen $k$; we take a \emph{single} ascent step \eqref{eq:lambda-update} per reverse step, so $\lambda_k$ tracks a slowly-moving target as the kernel $p_\theta(\cdot\mid x_{t_k})$ and the feasibility of the current sample evolve along the trajectory. This is a standard primal--dual / stochastic-approximation regime: the multiplier follows the constraint as decoding proceeds, tightening when the running constraint score lags the budget and relaxing once it is met. Proposition~\ref{prop:converge} describes the fixed-$k$ idealization that this online scheme tracks.
\end{remark}

\paragraph{Interpretation: escaping the fixed-$\gamma$ dilemma.} A constant $\gamma$ commits to one multiplier for the whole trajectory. Early in decoding, when $x_{t_k}$ is mostly masked and the classifier is uninformative, a large $\gamma$ over-tilts and pushes mass toward degenerate (e.g.\ invalid) regions; late in decoding, when feasibility is nearly settled, a small $\gamma$ under-tilts and the constraint is missed. Adaptive Dual Guidance instead applies \emph{just enough} multiplier to meet the budget at each step: $\lambda_k$ rises only while the constraint is unmet and decays otherwise, so the cumulative KL distortion to the base model is minimized subject to feasibility. This is precisely the behavior the experiments below exhibit.

% ======================================================================
\clearpage
\section{Experiments}
\label{sec:experiments}

\begin{table}[H]
\centering
\caption{Robustness across sampling budgets $K\in\{64,128,256,512\}$. SA constraint (QM9, MDLM, $N{=}500$, $\tau_{\text{train}}{=}3.0$); all rows within a block share the same time-conditioned discriminator, so the comparison is controlled. Fixed $\gamma\in\{0,1,3,5\}$ ($\gamma{=}0$ is the unguided base model) vs.\ the lowest-violation Adaptive Dual operating point of each block. Rates are N-independent; ``Yield'' is the count of valid, unique, novel \& SA-satisfying molecules. Arrows indicate whether higher ($\uparrow$) or lower ($\downarrow$) is better. ``Time'' is the wall-clock sampling time (seconds) for the $N{=}500$ samples on one GPU. Bold = lowest $\mathrm{Viol}_{3.0}$ within each block.}
\label{tab:steps}
\begin{tabular}{cl rrrrrr}
\toprule
$K$ & Method & Valid $\uparrow$ & Unique $\uparrow$ & $\mathrm{Viol}_{3.0}$ $\downarrow$ & Yield $\uparrow$ & QED $\uparrow$ & Time (s) $\downarrow$ \\
\midrule
\multirow{5}{*}{64}
  & Base ($\gamma{=}0$) & 59.6\% & 58.8\% & 84.90\% & 10 & 0.443 & 549 \\
  & $\gamma{=}1$        & 64.8\% & 51.2\% & 25.62\% & 46 & 0.493 & 543 \\
  & $\gamma{=}3$        & 65.0\% & 25.8\% & 12.00\% & 31 & 0.422 & 532 \\
  & $\gamma{=}5$        & 13.2\% & 4.6\%  & 34.85\% & 5  & 0.399 & 468 \\
  & Adaptive Dual       & 79.2\% & 40.8\% & \textbf{10.10\%} & 45 & 0.444 & 576 \\
\midrule
\multirow{5}{*}{128}
  & Base ($\gamma{=}0$) & 62.2\% & 60.8\% & 86.50\% & 8  & 0.468 & 794 \\
  & $\gamma{=}1$        & 70.8\% & 57.6\% & 21.47\% & 51 & 0.498 & 788 \\
  & $\gamma{=}3$        & 78.0\% & 31.6\% & 8.21\%  & 43 & 0.444 & 767 \\
  & $\gamma{=}5$        & 31.0\% & 8.6\%  & 23.23\% & 17 & 0.417 & 660 \\
  & Adaptive Dual       & 82.0\% & 47.0\% & \textbf{7.32\%}  & 52 & 0.482 & 845 \\
\midrule
\multirow{5}{*}{256}
  & Base ($\gamma{=}0$) & 64.2\% & 61.0\% & 85.05\% & 11 & 0.517 & 988 \\
  & $\gamma{=}1$        & 71.8\% & 57.8\% & 20.89\% & 55 & 0.510 & 977 \\
  & $\gamma{=}3$        & 80.6\% & 36.4\% & 6.95\%  & 46 & 0.459 & 955 \\
  & $\gamma{=}5$        & 54.2\% & 11.8\% & 9.96\%  & 25 & 0.437 & 826 \\
  & Adaptive Dual       & 76.0\% & 22.6\% & \textbf{6.05\%}  & 39 & 0.444 & 1044 \\
\midrule
\multirow{5}{*}{512}
  & Base ($\gamma{=}0$) & 64.0\% & 61.0\% & 86.25\% & 7  & 0.498 & 1132 \\
  & $\gamma{=}1$        & 66.0\% & 53.8\% & 19.70\% & 40 & 0.498 & 1116 \\
  & $\gamma{=}3$        & 86.0\% & 45.4\% & 4.42\%  & 51 & 0.460 & 1111 \\
  & $\gamma{=}5$        & 71.6\% & 15.2\% & 5.59\%  & 45 & 0.430 & 984 \\
  & Adaptive Dual       & 88.4\% & 43.0\% & \textbf{4.07\%}  & 44 & 0.492 & 1384 \\
\bottomrule
\end{tabular}
\end{table}

\begin{table}[H]
\centering
\caption{CDD-ALM (augmented-Lagrangian constrained discrete diffusion) on the SA constraint (QM9, MDLM, $\tau_{\text{train}}{=}3.0$), swept over inner/outer iterations and dual step $\eta$. Same metrics as Table~\ref{tab:steps}; here ``Yield'' is the novelty rate (\%). Arrows indicate whether higher ($\uparrow$) or lower ($\downarrow$) is better. ``Time'' is wall-clock sampling time (seconds). Bold = lowest $\mathrm{Viol}_{3.0}$.}
\label{tab:cdd-alm}
\begin{tabular}{l rrrrrr}
\toprule
Config (inner/outer/$\eta$) & Valid $\uparrow$ & Unique $\uparrow$ & $\mathrm{Viol}_{3.0}$ $\downarrow$ & Yield $\uparrow$ & QED $\uparrow$ & Time (s) $\downarrow$ \\
\midrule
5/20/0.05    & 25.0\% & 96.9\%  & 28.1\% & 28.1\% & 0.489 & 1144 \\
20/5/0.05    & 10.2\% & 100.0\% & 30.8\% & 30.8\% & 0.474 & 1562 \\
50/10/0.05   & 17.2\% & 100.0\% & 18.2\% & 27.3\% & 0.555 & 5203 \\
10/50/0.05   & 20.3\% & 100.0\% & 11.5\% & 15.4\% & 0.515 & 1237 \\
10/1000/0.05 & 18.0\% & 100.0\% & 21.7\% & 26.1\% & 0.496 & 1240 \\
5/20/0.1     & 18.0\% & 95.7\%  & \textbf{4.3\%} & 13.0\% & 0.529 & 1178 \\
20/5/0.1     & 14.1\% & 94.4\%  & 22.2\% & 33.3\% & 0.520 & 1562 \\
50/10/0.1    & 11.7\% & 100.0\% & 33.3\% & 33.3\% & 0.494 & 6458 \\
10/50/0.1    & 17.2\% & 100.0\% & 13.6\% & 22.7\% & 0.523 & 1441 \\
10/1000/0.1  & 15.6\% & 100.0\% & 25.0\% & 25.0\% & 0.521 & 1552 \\
\bottomrule
\end{tabular}
\end{table}

% ======================================================================
\begin{thebibliography}{9}
\bibitem{austin2021d3pm}
J.~Austin, D.~Johnson, J.~Ho, D.~Tarlow, R.~van den Berg.
\newblock Structured Denoising Diffusion Models in Discrete State-Spaces.
\newblock \emph{NeurIPS}, 2021.
\bibitem{lou2024sedd}
A.~Lou, C.~Meng, S.~Ermon.
\newblock Discrete Diffusion Modeling by Estimating the Ratios of the Data Distribution.
\newblock \emph{ICML}, 2024.
\bibitem{sahoo2024mdlm}
S.~Sahoo et al.
\newblock Simple and Effective Masked Diffusion Language Models.
\newblock \emph{NeurIPS}, 2024.
\bibitem{dhariwal2021guidance}
P.~Dhariwal, A.~Nichol.
\newblock Diffusion Models Beat GANs on Image Synthesis.
\newblock \emph{NeurIPS}, 2021.
\bibitem{nisonoff2024discreteguidance}
H.~Nisonoff, J.~Xiong, S.~Allenspach, J.~Listgarten.
\newblock Unlocking Guidance for Discrete State-Space Diffusion and Flow Models.
\newblock \emph{arXiv:2406.01572}, 2024.
\bibitem{cardei2025cdd}
M.~Cardei et al.
\newblock Constrained Discrete Diffusion.
\newblock \emph{arXiv:2503.09790}, 2025.
\bibitem{wu2023practical}
L.~Wu, B.~Trippe, C.~Naesseth, D.~Blei, J.~P.~Cunningham.
\newblock Practical and Asymptotically Exact Conditional Sampling in Diffusion Models.
\newblock \emph{NeurIPS}, 2023.
\bibitem{zhao2024probabilistic}
S.~Zhao et al.
\newblock Probabilistic Inference in Language Models via Twisted Sequential Monte Carlo.
\newblock \emph{ICML}, 2024.
\end{thebibliography}

\end{document}
